\section{INTRODUCTION}\label{sec:Introduction}
% Importance of satellite system monitoring
\IEEEPARstart{S}{atellite} fault diagnosis is critical for mission success in space operations, 
where maintenance opportunities are virtually nonexistent once deployed. 
The Telemetry, Tracking, and Command (TT\&C) system provides telemetry streams (electrical, thermal and status data) that 
enable continuous spacecraft health monitoring~\cite{ganesan2021fault}. 
Data-driven approaches can extract fault features from these telemetry streams~\cite{xiao2020deep}, 
yet their efficacy remains constrained in TT\&C centers by several factors: 
fault occurrences are infrequent, resulting in highly imbalanced datasets with scarce representative samples, 
and models trained on historical telemetry data often exhibit poor generalization capabilities when confronted with novel fault patterns emerging from diverse platforms 
or varying operational conditions~\cite{xie2021leo,bieber2023generic}.

% Feasibility of aggregating data across multiple TT&C centers
Aggregating fault data across multiple TT\&C centers would ideally address data scarcity and expand fault pattern coverage~\cite{von2018data}. 
By pooling telemetry data, each center could access more fault samples, diverse fault manifestations across platforms, 
and novel fault types absent in its local dataset—all essential for improving generalization to unseen faults in future operations. 
However, such direct data aggregation faces significant barriers in practice due to stringent data privacy requirements in aerospace applications. 
Satellite telemetry often contains sensitive information related to national security and privacy interests, 
rendering cross-center data sharing impractical.

% Introduction to federated learning
Federated learning~\cite{mcmahan2017communication} provides a promising privacy preserving alternative that reconciles these competing demands. 
By enabling collaborative model training without raw data exchange, 
federated learning preserves data sovereignty while allowing TT\&C centers to collectively benefit from shared fault pattern knowledge. 
In this paper, we investigate federated learning for satellite fault diagnosis, 
addressing two unique challenges specific to this domain:

\begin{figure}[t]
	\centering
	\includegraphics[width=\columnwidth]{fig/System Model/system model.pdf}
	\caption{System model.}
	\label{fig:system_model}
\end{figure}

% Problems addressed in this paper
\textit{Communication-limited}: A subset of TT\&C centers can participate in only a very small number of federated update rounds 
— in extreme cases only a single round — due to limited link availability, or operational constraints. 

\textit{Data Heterogeneity}: Telemetry data exhibits pronounced heterogeneity in two dimensions: 
(1) \textbf{feature-level}—identical faults manifest differently across centers due to platform and sensor diversity; 
(2) \textbf{label-level}—centers maintain disparate fault inventories with imbalanced distributions.

These challenges are illustrated in Fig.~\ref{fig:system_model}, where a communication-constrained TT\&C center joins at round 100 with new fault categories. 
Traditional federated learning then encounters: 
(1) \textbf{Model instability}—the model has difficulty retaining knowledge of old classes when learning new ones; 
(2) \textbf{Recency bias}—after the constrained center departs, the model gradually forgets its unique fault classes, as shown in Fig.~\ref{fig:problem-formulation}.

\begin{figure}[t]
	\centering
	\includegraphics[width=\columnwidth]{fig/Introduction/problem-formulation-acc.pdf}
  \caption{Problem formulation. $\alpha$ represents the degree of data heterogeneity.}
	\label{fig:problem-formulation}
\end{figure}

% Objective of this paper
The objective of this paper is to develop a model capable of accurately diagnosing faults across all TT\&C centers 
while minimizing the communication rounds needed for communication-constrained TT\&C centers. 
Related research on continual learning, federated learning with privacy concerns, and heterogeneous data handling is discussed in Section~\ref{sec:Related_Work}.

% Main contributions of this paper
In this study, we present a streaming federated learning framework for satellite fault diagnosis 
that ensures efficient communication and maintains data privacy. 
Under this framework, the TT\&C center with communication constraints only needs to 
participate in a \textbf{single round} of federated learning for the system to obtain a well-performing and stable model. 
% Summary of contributions
The key contributions of this work are outlined below:
\begin{itemize}
  \item 
  We provide a rigorous convergence analysis for the proposed algorithms in the communication-constrained satellite TT\&C setting. 
  Specifically, for the proposed algorithms under non-convex objectives, 
  we establish sublinear convergence to a neighborhood of a stationary point, 
  with a bias-controlled asymptotic error floor.
  \item 
  We design a streaming federated learning framework with two core algorithms: 
  (1) an Incremental Learning Algorithm based on Dual-Domain Knowledge Distillation for effective knowledge acquisition; 
  (2) a Federated Learning Algorithm based on Partial Forgetting Compensation to mitigate knowledge loss.
  \item 
  We perform extensive experiments, including comparisons with existing methods in the literature such as DTST and in-depth ablation studies. 
  Results show the framework's superior robustness under communication constraints and data heterogeneity.
\end{itemize}

% Paper organization
The remainder of this paper is structured as follows: 
Section \ref{sec:Related_Work} reviews related work on continual learning and federated learning. 
In Section \ref{sec:proposed_framework}, the proposed framework is detailed,
including the implementation procedures and the two core algorithms involved in this framework. 
Section \ref{sec:convergence-analysis} provides theoretical convergence analysis of the proposed algorithm.
Section \ref{sec:simulation-results} gives a series of case studies to demonstrate the effectiveness of the developed schemes. 
Section \ref{sec:conclusion} concludes the paper.